% Generated draft: requires mathematical and source review.
\documentclass[11pt]{article}
\usepackage[margin=25mm]{geometry}
\usepackage{fontspec}
\setmainfont{Latin Modern Roman}
\usepackage{amsmath,amssymb,mathtools,unicode-math}
\setmathfont{Latin Modern Math}
\usepackage{enumitem,needspace,xurl,hyperref}
\setlength{\parindent}{0pt}
\setlength{\parskip}{5pt}
\setlength{\emergencystretch}{3em}
\newcommand{\N}{\mathbb{N}}
\newcommand{\Z}{\mathbb{Z}}
\newcommand{\Q}{\mathbb{Q}}
\newcommand{\R}{\mathbb{R}}
\newcommand{\C}{\mathbb{C}}
\newcommand{\F}{\mathbb{F}}
\newcommand{\A}{\mathbb{A}}
\newcommand{\PP}{\mathbb{P}}
\newcommand{\E}{\mathbb{E}}
\newcommand{\eps}{\varepsilon}
\newcommand{\dd}{\,\mathrm{d}}
\begin{document}
\section*{UnsolvedMath ID: 800907}
\textbf{Source problem code:} catalogue entry
\subsection{Definitions and mathematical statement}
The supplied question does not specify the ambient space, the coefficient \(f\), the volume form denoted by \(dV\), the normalization of \(dd^c\), or the operator \(P\). These data must be fixed before a fully precise formulation is possible. In the equation below, \(V\) denotes the unknown function, whereas \(dV\) retains the source's notation for the volume element.

\textbf{Question 38 (attributed in the source to Eyssidieux--Guedj--Zeriahi).} Let \(v\) be a continuous viscosity supersolution of
\[
-(dd^c V)^n+f\,dV=0.
\]
Is \(P(v)\) a pluripotential supersolution of the same equation?

For clarity, a possible envelope interpretation can be expressed conditionally. On a specified complex domain or complex manifold \(X\) of complex dimension \(n\), let \(\operatorname{PSH}(X)\) denote the plurisubharmonic functions. If the intended operator is the plurisubharmonic envelope, its usual definition is
\[
P_X(v):=\left(\sup\{u\in\operatorname{PSH}(X):u\leq v\text{ on }X\}\right)^*,
\]
where \(h^*(x)=\limsup_{y\to x}h(y)\) denotes upper semicontinuous regularization. This identification of \(P\) is not specified in the supplied record. If it is intended, and if \(P_X(v)\) is locally bounded and the data define the relevant measures, the usual Bedford--Taylor interpretation of the requested conclusion is
\[
(dd^c P_X(v))^n\leq f\,dV
\]
as measures on \(X\). The source does not state hypotheses guaranteeing that this envelope or this measure interpretation is applicable, and it does not specify the viscosity supersolution convention or any boundary conditions.
\subsection{Short English statement}
The source asks whether applying \(P\) to a continuous viscosity supersolution of a complex Monge--Amp\`ere equation produces a pluripotential supersolution. The ambient space, equation data, and meaning of \(P\) remain unspecified; the envelope and measure formulation above is explicitly conditional.
\subsection{Sources}
\begin{itemize}
\item Local UnsolvedMath registry, permanent ID \(800907\), exact source problem code \texttt{catalogue entry}, titled \emph{Open Problems in Pluripotential Theory --- Question 38}. The registry supplies the status \texttt{open}; this is a source label, not an independently verified literature assessment.
\item Dinew, Guedj, Zeriahi, \emph{Open Problems in Pluripotential Theory} (2015), Question 38. Supplied source URL: \url{https://arxiv.org/abs/1511.00705}. The question itself carries the attribution Eyssidieux--Guedj--Zeriahi.
\item Supplied provenance records extraction by \texttt{arxiv\_latex} and access on \texttt{2026-10-16}. This access date is supplied metadata; no independent source visit or verification is asserted here.
\end{itemize}
\end{document}
